\section{Related Work}
\label{sec:related_work}

Stone's \textit{Pixel Perfect Timing Attacks with HTML5}~\cite{stone2013pixel} demonstrated the feasibility of cross-origin pixel stealing through browser rendering behavior. By exploiting software timing differences in the HTML5 \texttt{requestAnimationFrame} API, the attack measures how quickly CSS and SVG filters process pixels within a cross-origin iframe, allowing an attacker to infer image contents. This work established browser rendering operations as an unintended information channel and laid the foundation for subsequent SVG-based timing attacks. The following works all build upon this paper.

Andrysco et al.~\cite{andrysco2015subnormal} extended this approach by identifying hardware-level timing differences in floating-point operations. They demonstrated that arithmetic involving subnormal floating-point numbers can execute significantly slower than operations on normal numbers. By constructing SVG filters that trigger these operations, attackers can distinguish pixel values based on execution time. Unlike Stone's approach, which exploits browser rendering behavior, this attack leverages CPU-level floating-point performance characteristics, demonstrating how hardware behavior can undermine seemingly constant-time software.

Wang et al.'s \textit{GPU.zip}~\cite{wang2024gpuzip} expanded side-channel research to integrated GPU architectures by exploiting hardware-based graphical data compression. The authors demonstrated that undocumented, lossless compression mechanisms can introduce data-dependent differences in DRAM traffic and rendering time. By constructing carefully designed SVG filters that produce highly compressible pixel patterns, attackers can infer underlying pixel values. This work showed that hardware optimizations intended to improve performance can themselves introduce information leakage, extending earlier attacks beyond CPU execution timing.

O'Connell et al.'s \textit{Pixel Thief}~\cite{oconnell2024pixelthief} further advanced pixel-stealing attacks by exploiting CPU cache-state changes rather than relying directly on elapsed execution time. By monitoring the cache footprint left by SVG filter execution, the attack extracts pixel information despite modern browser timing mitigations, including reduced timer resolution and jitter.

Collectively, these works demonstrate the evolution of browser-based side-channel attacks from software-level rendering measurements to hardware-level information leakage. Stone established the feasibility of timing-based pixel extraction, Andrysco et al. identified CPU floating-point timing differences, GPU.zip exposed leakage through GPU memory compression, and Pixel Thief demonstrated cache-based extraction. Together, they highlight the limitations of software-level timing mitigations and the broader security implications of data-dependent hardware behavior.